\section{Problem Setup}\label{sec:setup}

\paragraph{Goal.}
Given a batch of trajectories with heterogeneous lengths and within-trajectory token dependence, we seek an unbiased linear aggregate of trajectory-level gradients that is statistically efficient. We cast the aggregation problem as estimation of a common mean under heteroscedastic noise whose variance scales with trajectory length and dependence structure, setting up a Bayesian treatment that will later recover and generalize the $\Delta L$ family of weights \citep{he2025deltal}.

\subsection{Setup and Notation}\label{subsec:setup-notation}

We consider an update step that aggregates a collection of $G\ge 1$ rollouts (trajectories) indexed by $i\in\{1,\dots,G\}$.
Trajectory $i$ has a (random) token length $L_i\in\mathbb{N}$ and yields a sequence of token-level gradient contributions
$\{\boldsymbol{Z}_{i,t}\}_{t=1}^{L_i}\subset\mathbb{R}^{p}$, where $p$ is the parameter dimension.
The (un-normalized) trajectory gradient is the token-sum
\begin{equation}
  \boldsymbol{g}_i \;:=\; \sum_{t=1}^{L_i}\boldsymbol{Z}_{i,t}.
  \label{eq:def_traj_grad_sum}
\end{equation}
Throughout, we regard the policy-gradient signal as a common mean vector $\boldsymbol{\mu}:=\mathbb{E}[\boldsymbol{g}_i]$ (under the current policy and sampling scheme),
and we study linear aggregators of the form \eqref{eq:linear-agg}.
The scaling constant $M>0$ in \eqref{eq:linear-agg} is fixed by the training objective (e.g., $M$ may correspond to a number of prompt groups in a batch, or to a normalization convention in the policy-gradient estimator).

To quantify \emph{dispersion} of a vector-valued gradient, we introduce a scalar functional
$\Psi:\mathbb{S}_{+}^{p}\to(0,\infty)$ on positive-semidefinite matrices (e.g., $\Psi(\boldsymbol{\Sigma})=\mathrm{tr}(\boldsymbol{\Sigma})$ or $\Psi(\boldsymbol{\Sigma})=\bm{u}^\top\boldsymbol{\Sigma}\bm{u}$ for a fixed direction $\bm{u}$)
and define the proxy variance
\begin{equation}
  v_i \;:=\; \Psi\!\left(\mathrm{Cov}(\boldsymbol{g}_i)\right).
  \label{eq:def_vi}
\end{equation}
The role of CAPO is to estimate how $v_i$ depends on length and within-trajectory dependence, and to use (an estimate of) $v_i^{-1}$ as the statistically efficient weighting rule.

\begin{table}[t]
\centering
\small
\begin{tabular}{ll}
\toprule
Symbol & Meaning \\
\midrule
$G$ & number of trajectories aggregated in the update \\
$L_i$ & token length of trajectory $i$ \\
$p$ & parameter dimension; $\boldsymbol{Z}_{i,t}\in\mathbb{R}^{p}$ \\
$\boldsymbol{Z}_{i,t}$ & token-level gradient contribution at position $t$ in trajectory $i$ \\
$\boldsymbol{g}_i$ & trajectory-level gradient $\sum_{t=1}^{L_i}\boldsymbol{Z}_{i,t}$ \\
$\boldsymbol{\mu}$ & common mean gradient, $\mathbb{E}[\boldsymbol{g}_i]=\boldsymbol{\mu}$ \\
$\widehat{\boldsymbol{g}}$ & aggregated gradient $\sum_{i=1}^{G} x_i\boldsymbol{g}_i$ \\
$x_i$ & trajectory weights; unbiasedness requires $\sum_i x_i = 1/M$ \\
$M$ & fixed scaling constant in \eqref{eq:linear-agg} \\
$\Psi(\cdot)$ & scalar functional on covariances used to define $v_i$ \\
$v_i$ & proxy dispersion $\Psi(\mathrm{Cov}(\boldsymbol{g}_i))$ \\
$\boldsymbol{\Gamma}_i(k)$ & token-level lag-$k$ autocovariance, \eqref{eq:Gamma_def} \\
$s(L;\xi)$ & dependence shape factor linking autocovariance to $v_i$ \\
$\beta$ & length exponent in the variance law $v_i\propto L_i^{\beta}$ \\
\bottomrule
\end{tabular}
\caption{Core notation used throughout the paper.}
\label{tab:notation}
\end{table}

\begin{assumption}[Cross-trajectory uncorrelatedness]\label{assump:cross_traj_uncorr}
For the trajectories aggregated in a single update, we assume
$\mathrm{Cov}(\boldsymbol{g}_i,\boldsymbol{g}_j)=\boldsymbol{0}$ for all $i\neq j$.
This assumption is exact when trajectories are sampled independently and is a standard working approximation when any remaining dependence is negligible compared to within-trajectory autocorrelation.
\end{assumption}

\begin{assumption}[Second-order stationarity within trajectories]\label{assump:stationary_tokens}
For each trajectory $i$, conditional on its length $L_i$, the token-level process $\{\boldsymbol{Z}_{i,t}\}_{t=1}^{L_i}$ is second-order stationary in the sense that
$\mathbb{E}[\boldsymbol{Z}_{i,t}]$ does not depend on $t$ and the lag-$k$ autocovariance
$\mathrm{Cov}(\boldsymbol{Z}_{i,t},\boldsymbol{Z}_{i,t-k})$ depends on $k$ but not on $t$ whenever both indices are in $\{1,\dots,L_i\}$.
\end{assumption}

Assumptions~\ref{assump:cross_traj_uncorr}--\ref{assump:stationary_tokens} are sufficient for the variance decompositions used in \S\ref{sec:method} and the appendices.
We discuss several dependence models for $\boldsymbol{\Gamma}_i(k)$ and their implications for the length-variance law in \S\ref{subsec:token_dependence} and \S\ref{subsec:kband-eta}.

\subsection{Observed Entities and Target Quantity}
We observe $G$ trajectories indexed by $i\in\{1,\dots,G\}$. Trajectory $i$ has length $L_i\in\mathbb{N}$ tokens and produces the un-normalized trajectory-level gradient.
\begin{equation}
  \boldsymbol{g}_i \;=\; \sum_{t=1}^{L_i} \boldsymbol{Z}_{i,t}\ \in\ \mathbb{R}^{p},
  \qquad
  \mathbb{E}\!\left[\,\boldsymbol{g}_i\,\right] \;=\; \boldsymbol{\mu}\ \in\ \mathbb{R}^{p},
  \label{eq:traj-grad}
\end{equation}
where $\boldsymbol{Z}_{i,t}$ is the token-level score contribution (e.g., advantage-weighted policy-score increment). The learning signal $\boldsymbol{\mu}$ is common across trajectories in a batch.

\subsection{A Statistically Efficient, Linear, Unbiased Estimator}
Training aggregates the per-trajectory gradients linearly,
\begin{equation}
  \widehat{\boldsymbol{g}}
  \;=\;
  \sum_{i=1}^{G} x_i\,\boldsymbol{g}_i,
  \qquad
  \sum_{i=1}^{G} x_i \;=\; \frac{1}{M},
  \label{eq:linear-agg}
\end{equation}
for a fixed scaling constant $M>0$. Under \eqref{eq:traj-grad}, any choice of weights satisfying the normalization in \eqref{eq:linear-agg} yields an unbiased estimate, $\mathbb{E}[\widehat{\boldsymbol{g}}]=\boldsymbol{\mu}/M$. The design problem is to choose $\{x_i\}$ to minimize the dispersion of $\widehat{\boldsymbol{g}}$ subject to unbiasedness.

\subsection{Non-iid Tokens}\label{subsec:token_dependence}
Token contributions within a trajectory are dependent. Define the auto-covariance function,
\begin{equation}
  \boldsymbol{\Gamma}_i(k)
  \;:=\;
  \mathrm{Cov}\!\big(\boldsymbol{Z}_{i,t},\,\boldsymbol{Z}_{i,t-k}\big)\in\mathbb{R}^{p\times p},
  \qquad k\in\mathbb{Z}.
  \label{eq:Gamma_def}
\end{equation}
Then the trajectory covariance, under stationarity condition, consult Appendix \ref{app:bartlett}, is
\begin{equation}
  \mathrm{Var}(\boldsymbol{g}_i)
  \;=\;
  \sum_{k=-(L_i-1)}^{L_i-1} (L_i-|k|)\,\boldsymbol{\Gamma}_i(k).
  \label{eq:var-sum}
\end{equation}
For aggregation we summarize $\mathrm{Var}(\boldsymbol{g}_i)$ by a positive scalar functional
\begin{equation}
  v_i \;:=\; \Psi\!\big(\mathrm{Var}(\boldsymbol{g}_i)\big)\ \in\ (0,\infty),
  \qquad
  \text{e.g.,}\ \Psi(\boldsymbol{A})=\mathrm{tr}(\boldsymbol{A})\ \text{or}\ \Psi(\boldsymbol{A})=\mathbb{E}_{\boldsymbol{u}}[\boldsymbol{u}^\top \boldsymbol{A}\boldsymbol{u}],
  \label{eq:scalar-proxy}
\end{equation}
where $\boldsymbol{u}$ is a unit vector (fixed or randomized). If $\{v_i\}$ were known, the minimum-variance unbiased (MVU) linear weights are
\begin{equation}
  x_i^\star \;=\; \frac{1}{M}\,\frac{v_i^{-1}}{\sum_{j=1}^{G} v_j^{-1}}.
  \label{eq:mvu-weights}
\end{equation}
For a detailed proof along with the assumptions consult Appendix \ref{app:mvu-weights}. Hence the core statistical task is to model/estimate $v_i$ as a function of length and dependence.


% \section{Problem Setup}\label{sec:setup}

% \paragraph{Goal.}
% Given a batch of trajectories with heterogeneous lengths and within-trajectory token dependence, we seek an unbiased linear aggregate of trajectory-level gradients that is statistically efficient. We cast the aggregation problem as estimation of a common mean under heteroscedastic noise whose variance scales with trajectory length and dependence structure, setting up a Bayesian treatment that will later recover and generalize the $\Delta L$ family of weights.

% \subsection{Observed Entities and Target Quantity}
% We observe $G$ trajectories (responses) indexed by $i\in\{1,\dots,G\}$. Trajectory $i$ has length $L_i\in\mathbb{N}$ tokens and produces the (unnormalized) trajectory-level gradient
% \begin{equation}
%   \boldsymbol{g}_i \;=\; \sum_{t=1}^{L_i} \boldsymbol{Z}_{i,t}\ \in\ \mathbb{R}^{p},
%   \qquad
%   \mathbb{E}\!\left[\,\boldsymbol{g}_i\,\right] \;=\; \boldsymbol{\mu}\ \in\ \mathbb{R}^{p},
%   \label{eq:traj-grad}
% \end{equation}
% where $\boldsymbol{Z}_{i,t}$ is the token-level score contribution (e.g., advantage-weighted policy-score increment). The learning signal $\boldsymbol{\mu}$ is common across trajectories in a batch.

% \subsection{Design Objective: Linear, Unbiased, and Efficient}
% Training aggregates the per-trajectory gradients linearly,
% \begin{equation}
%   \widehat{\boldsymbol{g}}
%   \;=\;
%   \sum_{i=1}^{G} x_i\,\boldsymbol{g}_i,
%   \qquad
%   \sum_{i=1}^{G} x_i \;=\; \frac{1}{M},
%   \label{eq:linear-agg}
% \end{equation}
% for a fixed scaling constant $M>0$. Under \eqref{eq:traj-grad}, any choice of weights satisfying \eqref{eq:linear-agg} yields an unbiased estimate, $\mathbb{E}[\widehat{\boldsymbol{g}}]=\boldsymbol{\mu}/M$. The design problem is to choose $\{x_i\}$ to minimize the dispersion of $\widehat{\boldsymbol{g}}$ subject to unbiasedness.

% \subsection{Non-iid Tokens and a Scalar Variance Proxy}
% Token contributions within a trajectory are dependent. Define the auto-covariance function (along the auto-regressive index)
% \begin{equation}
%   \boldsymbol{\Gamma}_i(k)
%   \;:=\;
%   \mathrm{Cov}\!\big(\boldsymbol{Z}_{i,t},\,\boldsymbol{Z}_{i,t-k}\big)\in\mathbb{R}^{p\times p},
%   \qquad k\in\mathbb{Z}.
% \end{equation}
% Then the trajectory covariance, under stationarity, is
% \begin{equation}
%   \mathrm{Var}(\boldsymbol{g}_i)
%   \;=\;
%   \sum_{k=-(L_i-1)}^{L_i-1} (L_i-|k|)\,\boldsymbol{\Gamma}_i(k).
%   \label{eq:var-sum}
% \end{equation}
% For aggregation we summarize $\mathrm{Var}(\boldsymbol{g}_i)$ by a positive scalar functional
% \begin{equation}
%   v_i \;:=\; \Psi\!\big(\mathrm{Var}(\boldsymbol{g}_i)\big)\ \in\ (0,\infty),
%   \qquad
%   \text{e.g.,}\ \Psi(\boldsymbol{A})=\mathrm{tr}(\boldsymbol{A})\ \text{or}\ \Psi(\boldsymbol{A})=\mathbb{E}_{\boldsymbol{u}}[\boldsymbol{u}^\top \boldsymbol{A}\boldsymbol{u}],
%   \label{eq:scalar-proxy}
% \end{equation}
% where $\boldsymbol{u}$ is a unit vector (fixed or randomized). If $\{v_i\}$ were known, the minimum-variance unbiased (MVU) linear weights are
% \begin{equation}
%   x_i^\star \;=\; \frac{1}{M}\,\frac{v_i^{-1}}{\sum_{j=1}^{G} v_j^{-1}}.
%   \label{eq:mvu-weights}
% \end{equation}
% Hence the core statistical task is to model/estimate $v_i$ as a function of length and dependence.





% \section{Problem Setup}

% \paragraph{Goal.}
% Given a batch of trajectories with heterogeneous lengths and within-trajectory token dependence, we seek an unbiased linear aggregate of trajectory-level gradients that is statistically efficient. We cast the aggregation problem as estimation of a common mean under heteroscedastic noise whose variance scales with trajectory length and dependence structure, setting up a Bayesian treatment that will later recover and generalize the $\Delta L$ family of weights.

% \subsection{Observed Entities and Target Quantity}
% We observe $G$ trajectories (responses) indexed by $i\in\{1,\dots,G\}$. Trajectory $i$ has length $L_i\in\mathbb{N}$ tokens and produces the (unnormalized) trajectory-level gradient
% \begin{equation}
%   \boldsymbol{g}_i \;=\; \sum_{t=1}^{L_i} \boldsymbol{Z}_{i,t}\ \in\ \mathbb{R}^{p},
%   \qquad
%   \mathbb{E}\!\left[\,\boldsymbol{g}_i\,\right] \;=\; \boldsymbol{\mu}\ \in\ \mathbb{R}^{p},
%   \label{eq:traj-grad}
% \end{equation}
% where $\boldsymbol{Z}_{i,t}$ is the token-level score contribution (e.g., advantage-weighted policy-score increment). The learning signal $\boldsymbol{\mu}$ is common across trajectories in a batch.

% \subsection{A Statistically Efficient, Linear, Unbiased Estimator}
% Training aggregates the per-trajectory gradients linearly,
% \begin{equation}
%   \widehat{\boldsymbol{g}}
%   \;=\;
%   \sum_{i=1}^{G} x_i\,\boldsymbol{g}_i,
%   \qquad
%   \sum_{i=1}^{G} x_i \;=\; \frac{1}{M},
%   \label{eq:linear-agg}
% \end{equation}
% for a fixed scaling constant $M>0$. Under \eqref{eq:traj-grad}, any choice of weights satisfying the normalization in \eqref{eq:linear-agg} yields an unbiased estimate, $\mathbb{E}[\widehat{\boldsymbol{g}}]=\boldsymbol{\mu}/M$. The design problem is to choose $\{x_i\}$ to minimize the dispersion of $\widehat{\boldsymbol{g}}$ subject to unbiasedness.

% \subsection{Non-iid Tokens}
% Token contributions within a trajectory are dependent. Define the auto-covariance function (along the auto-regressive index)
% \begin{equation}
%   \boldsymbol{\Gamma}_i(k)
%   \;:=\;
%   \mathrm{Cov}\!\big(\boldsymbol{Z}_{i,t},\,\boldsymbol{Z}_{i,t-k}\big)\in\mathbb{R}^{p\times p},
%   \qquad k\in\mathbb{Z}.
% \end{equation}
% Then the trajectory covariance, under stationarity condition, is
% \begin{equation}
%   \mathrm{Var}(\boldsymbol{g}_i)
%   \;=\;
%   \sum_{k=-(L_i-1)}^{L_i-1} (L_i-|k|)\,\boldsymbol{\Gamma}_i(k).
%   \label{eq:var-sum}
% \end{equation}
% For aggregation we summarize $\mathrm{Var}(\boldsymbol{g}_i)$ by a positive scalar functional
% \begin{equation}
%   v_i \;:=\; \Psi\!\big(\mathrm{Var}(\boldsymbol{g}_i)\big)\ \in\ (0,\infty),
%   \qquad
%   \text{e.g.,}\ \Psi(\boldsymbol{A})=\mathrm{tr}(\boldsymbol{A})\ \text{or}\ \Psi(\boldsymbol{A})=\mathbb{E}_{\boldsymbol{u}}[\boldsymbol{u}^\top \boldsymbol{A}\boldsymbol{u}],
%   \label{eq:scalar-proxy}
% \end{equation}
% where $\boldsymbol{u}$ is a unit vector (fixed or randomized). If $\{v_i\}$ were known, the minimum-variance unbiased (MVU) linear weights are
% \begin{equation}
%   x_i^\star \;=\; \frac{1}{M}\,\frac{v_i^{-1}}{\sum_{j=1}^{G} v_j^{-1}}.
%   \label{eq:mvu-weights}
% \end{equation}
% Hence the core statistical task is to model/estimate $v_i$ as a function of length and dependence.

% \subsection{Length-Only Benchmark and the $\Delta L$ Family}
% Under an i.i.d.\ token model with $\mathrm{Var}(\boldsymbol{Z}_{i,t})=\boldsymbol{\Sigma}_0$ and $\boldsymbol{\Gamma}_i(k)=\boldsymbol{0}$ for $k\neq 0$,
% \begin{equation}
%   \mathrm{Var}(\boldsymbol{g}_i)=L_i\,\boldsymbol{\Sigma}_0
%   \quad\Longrightarrow\quad
%   v_i \propto L_i.
% \end{equation}
% Substituting into \eqref{eq:mvu-weights} gives $x_i\propto L_i^{-1}$, i.e., the $\Delta L$ rule with exponent $\alpha=1$. More generally, the paper proposes a one-parameter family
% \begin{equation}
%   x_i \ \propto\ L_i^{-\alpha},\qquad \alpha\in[0,1],
% \end{equation}
% interpolating between constant weighting ($\alpha=0$) and inverse-length weighting ($\alpha=1$), as a practical control of variance induced by heterogeneous lengths.


% % === Replacement Section: Bayesian Setup under a General Prior, with Derivations ===

% \subsection{Bayesian Setup under a General Prior}
% We project each trajectory-level gradient onto a fixed direction to obtain a scalar summary. Let $\bm{u}$ be a unit vector and define
% \begin{equation}
%   g_i \;=\; \bm{u}^\top \bm{g}_i,\qquad \mathbb{E}[g_i]=\mu,\qquad \mathrm{Var}(g_i)=v_i.
%   \label{eq:def_gi_scalar}
% \end{equation}
% Conditional on the variance parameters $\vartheta := (\sigma^2, \beta, \xi)$,

% \begin{equation}
%   g_i \,\big|\, \mu,\vartheta \ \sim\ \mathcal{N}\!\big(\mu,\ v_i(\vartheta)\big),
%   \qquad
%   v_i(\vartheta)\;=\;\sigma^2\,L_i^{\beta}\,s(L_i;\,\xi),
%   \label{eq:hetero_general_prior}
% \end{equation}

% where $\vartheta=(\sigma^2,\beta,\xi)$. Here $\sigma^2>0$ is a global scale, $\beta\in[0,2]$ controls length scaling, and $s(L_i;\xi)>0$ is a shape factor encoding within-trajectory dependence. This modular treatment allows us to plug in different estimates for the shape factor neatly.

% \paragraph{General priors.}
% We adopt a flat prior on the mean and a flexible, proper prior on the variance and shape parameters:
% \begin{equation}
%   p(\mu)\ \propto\ 1,\qquad
%   \sigma^2 \sim \mathrm{Inv\text{-}Gamma}(a_0,b_0),\qquad
%   \beta \sim \pi_\beta,\qquad
%   \xi \sim \pi_\xi.
%   \label{eq:priors_general}
% \end{equation}
% Jeffreys' prior for scale is recovered in the limit $a_0,b_0\to 0$. The prior $\pi_\beta$ can be chosen agnostic (centered at $1$ with moderate variance) or informative to reflect anticipated sub/superlinear growth. We later demonstrate the relationship between this and the 'correctness' of the trajectory. The prior $\pi_\xi$ encodes beliefs about token dependence.

% \paragraph{Likelihood geometry and sufficient sums.}
% Define precisions and aggregated quantities
% \begin{equation}
%   \lambda_i(\vartheta) := \frac{1}{v_i(\vartheta)} = \frac{1}{\sigma^2}\,\omega_i(\beta,\xi),
%   \qquad
%   \omega_i(\beta,\xi) := L_i^{-\beta}\,s(L_i;\xi)^{-1},
%   \label{eq:def_lambda_omega}
% \end{equation}
% \begin{equation}
%   \Lambda(\vartheta) \;=\; \sum_{i=1}^G \lambda_i(\vartheta),\qquad
%   T(\vartheta) \;=\; \sum_{i=1}^G \lambda_i(\vartheta)\,g_i.
%   \label{eq:Lambda_T}
% \end{equation}
% The conditional posterior for $\mu$ given $\vartheta$ is Normal:
% \begin{equation}
%   \mu\,\big|\,\vartheta,\{g_i\} \ \sim\ \mathcal{N}\!\big(m(\vartheta),\,\Lambda(\vartheta)^{-1}\big),
%   \qquad
%   m(\vartheta)=\frac{T(\vartheta)}{\Lambda(\vartheta)}=\frac{\sum_i \lambda_i(\vartheta) g_i}{\sum_i \lambda_i(\vartheta)}.
%   \label{eq:mu_conditional}
% \end{equation}
% Let $m(\beta,\xi)$ denote $m(\vartheta)$ with $\sigma^2$ profiled out via \eqref{eq:def_lambda_omega}, and define the weighted residual sum of squares
% \begin{equation}
%   \mathrm{RSS}_\omega(\beta,\xi) \;=\; \sum_{i=1}^G \omega_i(\beta,\xi)\,\big(g_i - m(\beta,\xi)\big)^2.
%   \label{eq:rss_omega}
% \end{equation}
% Combining \eqref{eq:priors_general}--\eqref{eq:rss_omega}, the conditional posterior of $\sigma^2$ is
% \begin{equation}
%   \sigma^2 \,\big|\, \beta,\xi,\{g_i\} \ \sim\ \mathrm{Inv\text{-}Gamma}\!\left(a_0+\frac{G-1}{2},\ b_0+\frac{1}{2}\,\mathrm{RSS}_\omega(\beta,\xi)\right).
%   \label{eq:sigma2_posterior}
% \end{equation}

% \paragraph{Bayes weights (posterior-precision aggregator).}
% Define
% \begin{equation}
%   \widetilde{w}_i \;:=\; \mathbb{E}\!\left[\lambda_i(\vartheta)\,\big|\,\{g_j\}\right]
%   \;=\; \mathbb{E}\!\left[\frac{1}{\sigma^2}\,\omega_i(\beta,\xi)\,\Big|\,\{g_j\}\right].
%   \label{eq:bayes_weight_def}
% \end{equation}
% A practical estimator uses either (i) a Laplace/MAP plug-in $(\hat\beta,\hat\xi)$ yielding $\widetilde{w}_i \propto \omega_i(\hat\beta,\hat\xi)$, or (ii) Monte Carlo averaging over draws $(\beta_s,\xi_s,\sigma^2_s)$ from the posterior. The Bayes-optimal unbiased linear aggregate then takes
% \begin{equation}
%   x_i^{\mathrm{Bayes}} \;=\; \frac{1}{M}\,\frac{\widetilde{w}_i}{\sum_{j=1}^G \widetilde{w}_j}.
%   \label{eq:bayes_weight_rule}
% \end{equation}

% \subsection{Jeffreys' Prior Baseline: Recovering $\Delta L$ with $\alpha=1$}
% Assume independence ($s\equiv 1$) and fix $\beta=1$. Let $a_0,b_0\to 0$ so that $p(\sigma^2)\propto 1/\sigma^2$. From \eqref{eq:def_lambda_omega}--\eqref{eq:sigma2_posterior},
% \begin{equation}
%   \lambda_i(\vartheta) \;=\; \frac{1}{\sigma^2}\,L_i^{-1},
%   \qquad
%   \mathbb{E}\!\left[\sigma^{-2}\,\big|\,\beta=1,\{g_i\}\right] \;=\; \frac{G-1}{\,\mathrm{RSS}_\omega(\beta{=}1)\,},
% \end{equation}
% so the common factor $\mathbb{E}[\sigma^{-2}\mid\cdot]$ cancels in \eqref{eq:bayes_weight_rule}, and
% \begin{equation}
%   x_i^{\mathrm{Bayes}} \ \propto\ L_i^{-1}.
%   \label{eq:deltaL_alpha1}
% \end{equation}
% Thus the $\Delta L$ rule with exponent $\alpha=1$ is exactly the Bayes posterior-precision aggregator under Jeffreys' prior and the i.i.d.\ length law.

% \subsection{Informative Prior on the Length Exponent: $\alpha<1$}
% Retain independence ($s\equiv 1$) but let $\beta$ be random with a proper prior $\pi_\beta$ and posterior $p(\beta\mid\{g_i\})$. Using \eqref{eq:bayes_weight_def},
% \begin{equation}
%   \widetilde{w}_i \;\propto\; \mathbb{E}\!\left[L_i^{-\beta}\,\big|\,\{g_j\}\right].
%   \label{eq:wtilde_beta_random}
% \end{equation}
% Two useful consequences:
% \begin{enumerate}
%   \item \emph{MAP/Laplace approximation.} If $p(\beta\mid\{g_i\})$ is sharply peaked at $\hat\beta$, then
%   \begin{equation}
%     \widetilde{w}_i \ \propto\ L_i^{-\hat\beta}
%     \quad\Longrightarrow\quad
%     x_i^{\mathrm{Bayes}} \ \propto\ L_i^{-\hat\beta},
%     \label{eq:laplace_alpha}
%   \end{equation}
%   i.e., the $\Delta L$ family with \(\alpha=\hat\beta<1\) when posterior mass for $\beta$ concentrates below $1$ (sublinear variance growth).
%   \item \emph{Closed-form moment under Gaussian posterior.} If $\beta\mid\{g_i\}\sim \mathcal{N}(m_\beta,\tau_\beta^2)$ (ignoring truncation),
%   \begin{equation}
%     \mathbb{E}\!\left[L_i^{-\beta}\mid\{g_i\}\right]
%     \;=\;
%     \exp\!\Big(-m_\beta\ln L_i + \tfrac{1}{2}\tau_\beta^2(\ln L_i)^2\Big)
%     \;\approx\; L_i^{-m_\beta}\ \text{ when }\ \tau_\beta^2\text{ is small},
%     \label{eq:moment_gaussian_beta}
%   \end{equation}
%   so an effective exponent is \(\alpha \approx m_\beta < 1\) when the posterior mean of $\beta$ is subunit.
% \end{enumerate}
% In both cases the Bayes rule coincides with $\Delta L$ but with a \emph{data-driven} exponent $\alpha$ induced by an informative prior on $\beta$. Dependence corrections $s(L_i;\xi)$ can be reintroduced by replacing $L_i^{-\beta}$ with $L_i^{-\beta}s(L_i;\xi)^{-1}$ in \eqref{eq:wtilde_beta_random}.


% \subsection{Why $\alpha<1$ When Longer Responses Are More Likely Correct}

% \paragraph{Intuition.}
% If longer responses have higher probability of being correct, then---after baselineing---their token contributions are less noisy on average. In a heteroskedastic model, this means the \emph{per-token variance} decreases with length, so the \emph{sum} over $L$ tokens grows \emph{sublinearly} in $L$. Consequently, inverse-variance weights scale as $L^{-\alpha}$ with $\alpha<1$.

% \paragraph{Heteroskedastic variance model tied to correctness.}
% Let $g_i=\sum_{t=1}^{L_i} Z_{i,t}$ be the scalarized trajectory gradient (Sec.~\S\ref{eq:traj-grad}–\S\ref{eq:hetero-model} in your setup). Suppose the per-token variance contracts with response length via a correctness factor $\kappa(L)\in(0,1]$:
% \begin{equation}
%   \mathrm{Var}(Z_{i,t}\mid L_i=L)\;=\;\sigma^2\,\kappa(L),
%   \qquad
%   \kappa'(L)\le 0.
%   \label{eq:per-token-variance}
% \end{equation}
% A concrete rationale is a Bernoulli-like reward component with success probability $p(L)$ (``response is correct'') and a baseline $b(L)\approx \mathbb{E}[r\mid L]=p(L)$, which yields a multiplicative factor $\kappa(L)\propto p(L)\,\bigl(1-p(L)\bigr)$ in the gradient noise. As $p(L)\uparrow 1$, we have $\kappa(L)\downarrow 0$.

% Accounting for token dependence via $s(L;\xi)$ (Sec.~\S\ref{eq:hetero-model}), the trajectory variance takes the form
% \begin{equation}
%   \mathrm{Var}(g_i\mid L_i=L)\;=\;\sigma^2\,L\,\kappa(L)\,s(L;\xi).
%   \label{eq:traj-var-kappa}
% \end{equation}
% Assume a regularly varying contraction:
% \begin{equation}
%   \kappa(L)\;=\;c\,L^{-\delta}\,(1+o(1)),
%   \qquad
%   \delta\in[0,1),
%   \label{eq:kappa-regular}
% \end{equation}
% which is implied, for example, if $p(L)=1-\tilde c\,L^{-\gamma}+o(L^{-\gamma})$ with $\gamma>0$, since then
% \(
% p(L)\bigl(1-p(L)\bigr)
% =
% \tilde c\,L^{-\gamma}+o(L^{-\gamma})
% \)
% and one may set $\delta=\gamma$ up to constants.

% Combining \eqref{eq:traj-var-kappa}–\eqref{eq:kappa-regular} gives the effective length law
% \begin{equation}
%   v(L)\;:=\;\mathrm{Var}(g_i\mid L_i=L)
%   \;=\;\sigma^2\,L^{\,1-\delta}\,s(L;\xi)\,\bigl(1+o(1)\bigr).
%   \label{eq:effective-length-law}
% \end{equation}
% Thus the variance exponent is $\beta=1-\delta<1$ whenever correctness improves with length at a polynomial rate.

% \paragraph{Bayesian aggregation and the $\Delta L$ exponent.}
% Under the heteroskedastic Normal working model with
% \(
% v_i(\vartheta)=\sigma^2\,L_i^{\beta}\,s(L_i;\xi)
% \)
% and Jeffreys' prior on $\sigma^2$ (scale invariance), the Bayes posterior-precision weights satisfy, up to a common factor,
% \begin{equation}
%   x_i^{\mathrm{Bayes}}
%   \;\propto\;
%   \frac{1}{L_i^{\beta}\,s(L_i;\xi)}.
%   \label{eq:bayes-weight-beta}
% \end{equation}
% If $s(L;\xi)\equiv 1$ (independence benchmark), \eqref{eq:bayes-weight-beta} is exactly the $\Delta L$ family with exponent
% \(
% \alpha=\beta.
% \)
% By \eqref{eq:effective-length-law}, correctness-driven contraction implies $\beta=1-\delta<1$, hence
% \begin{equation}
%   x_i \;\propto\; L_i^{-\alpha},
%   \qquad
%   \alpha \;=\; 1-\delta \;<\; 1.
%   \label{eq:alpha-less-than-one}
% \end{equation}

% \paragraph{Examples linking $p(L)$ to $\alpha$.}
% \begin{enumerate}
% \item \emph{Power-law approach to correctness.}
% If $p(L)=1-\tilde c\,L^{-\gamma}+o(L^{-\gamma})$ with $\gamma\in(0,1)$, then
% \(
% \kappa(L)\propto L^{-\gamma}
% \)
% and
% \(
% \alpha=1-\gamma\in(0,1).
% \)
% \item \emph{Logistic saturation.}
% If $p(L)=\bigl(1+\exp\{-a(\log L-b)\}\bigr)^{-1}$ with $a>0$, then for large $L$,
% \(
% 1-p(L)\asymp L^{-a}
% \),
% so again $\alpha=1-a$ (clipped to $(0,1)$).
% \end{enumerate}

% \paragraph{With dependence $s(L;\xi)\neq 1$.}
% All conclusions persist with serial dependence: replace $L^{-\beta}$ by $L^{-\beta} s(L;\xi)^{-1}$ in \eqref{eq:bayes-weight-beta}. For AR(1), $s(L;\xi)$ approaches a positive constant as $L\to\infty$, so the large-$L$ exponent remains $\alpha=1-\delta$. For $k$-banded dependence, $s(L;\xi)=1+\tfrac{2}{L}\sum_{h=1}^{k}(L-h)\rho_h=1+O(1)$, again preserving $\alpha=1-\delta$ asymptotically.

% \paragraph{Informative priors strengthen $\alpha<1$.}
% If $\beta$ is random with a proper prior concentrated below $1$ (reflecting the belief ``longer answers are cleaner''), then
% \(
% x_i^{\mathrm{Bayes}}\propto \mathbb{E}[L_i^{-\beta}\mid\text{data}]
% \)
% and a Laplace/mean summary yields \(\alpha\simeq \mathbb{E}[\beta\mid\text{data}]<1\). This is the principled Bayesian route to $\alpha<1$ when correctness improves with length.


%%%%%% OLDEST FRAME

% \subsection{A Principled Bayesian Weighting Scheme}
% \label{subsec:bayesian}
% In order to put this into a unified Bayesian framework, we use a scalar projection of each trajectory gradient. Let $\boldsymbol{u}\in\mathbb{R}^{p}$ be fixed (e.g., an arbitrary unit vector), and define
% \begin{equation}
%   g_i \;:=\; \boldsymbol{u}^\top \boldsymbol{g}_i \in \mathbb{R},
%   \qquad
%   \mathbb{E}[g_i]=\mu \in \mathbb{R},
%   \qquad
%   \mathrm{Var}(g_i)=v_i.
% \end{equation}

% Given a normally distributed trajectory $g_i$ with hetero-scedastic variance, 
% \begin{equation}
%   g_i \,\big|\, \mu,\vartheta \ \sim\ \mathcal{N}\!\big(\mu,\ v_i(\vartheta)\big),
%   \qquad
%   v_i(\vartheta)\;=\;\sigma^2\,L_i^{\beta}\,s(L_i;\,\xi),
%   \label{eq:hetero-model}
% \end{equation}
% with variance parameters $\vartheta=(\sigma^2,\beta,\xi)$. Here $\sigma^2>0$ is a global scale, $\beta\in[0,2]$ controls length scaling (the i.i.d. benchmark is $\beta=1$), and $s(L_i;\xi)>0$ is a shape factor encoding within-trajectory dependence structure.

% \paragraph{Priors.}
% We employ a noninformative prior on the mean and a scale-invariant prior on the variance,
% \begin{equation}
%   p(\mu)\ \propto\ 1,
%   \qquad
%   p(\sigma^2)\ \propto\ \frac{1}{\sigma^2},
%   \qquad
%   \beta \sim \pi_\beta,\ \ \xi \sim \pi_\xi,
% \end{equation}
% with $\pi_\beta$ and $\pi_\xi$ specified to capture either agnostic or informed beliefs about length scaling and dependence.

% \subsection{Roadmap to Bayesian Identifications (Preview)}
% Under the model \eqref{eq:hetero-model}, the Bayes-optimal linear unbiased aggregate uses weights proportional to the posterior mean precision, i.e.,
% \begin{equation}
%   \widetilde{w}_i \;:=\; \mathbb{E}\!\left[\frac{1}{v_i(\vartheta)}\,\bigg|\,\{g_j\}_{j=1}^G\right]
%   \;=\;
%   \mathbb{E}\!\left[\frac{1}{\sigma^2\,L_i^{\beta}\,s(L_i;\xi)}\,\bigg|\,\{g_j\}\right],
%   \qquad
%   x_i^{\mathrm{Bayes}} \;=\; \frac{1}{M}\,\frac{\widetilde{w}_i}{\sum_{j=1}^{G}\widetilde{w}_j}.
%   \label{eq:bayes-weights}
% \end{equation}
% In the independence benchmark $s\equiv 1$ with a noninformative prior on $\sigma^2$ and $\beta=1$, \eqref{eq:bayes-weights} reduces to $x_i\propto L_i^{-1}$ (the $\Delta L$ rule with $\alpha=1$). Proper priors on $\beta$ yield posterior summaries that act as data-driven exponents ($\alpha\approx \mathbb{E}[\beta\mid\{g_j\}]$), while explicit dependence templates $s(L;\xi)$ (compound symmetry, $k$-banded, geometric) introduce principled corrections to the effective length law. These derivations follow in the next section.